\chapter{Raccolta Tematica di Esercizi ed Esami d'Appello Risolti}
\label{chap:raccolta_esami}

\section*{Presentazione della Raccolta}
Questa sezione raccoglie un compendio esaustivo e metodico di \textbf{temi d'esame ufficiali}, prove in itinere (compitini) ed esercizi d'appello del corso di \textbf{Analisi Matematica 1} (tenuto dal Prof. Giovanni Alberti, Università di Pisa), integrati con le relative \textbf{soluzioni dettagliate passo dopo passo}.

Gli esercizi sono stati accuratamente classificati e organizzati per \textbf{aree tematiche} per consentire una preparazione mirata e progressiva in vista delle prove scritta e orale:

\begin{enumerate}[label=\textbf{\Roman*.}]
    \item \textbf{Limiti, Sviluppi di Taylor, Simboli di Landau e Parti Principali}: calcolo di limiti mediante ordini di grandezza ($\ll$), cancellazioni asintotiche, sviluppi di Maclaurin e dipendenza da parametri reali $\alpha, a, \beta$.
    \item \textbf{Calcolo Differenziale, Punti Critici e Studio Qualitativo di Funzione}: ricerca di massimi e minimi assoluti su intervalli, calcolo delle rette tangenti, studio di concavità/convessità e risoluzione grafica di equazioni/disequazioni.
    \item \textbf{Calcolo Integrale, Primitive, Aree, Volumi di Rotazione e Integrali Impropri}: integrazione per parti e sostituzione, decomposizione in fratti semplici, calcolo di aree piane e volumi di solidi di rotazione attorno agli assi cartesiani o a rette generiche, e convergenza di integrali impropri con parametro.
    \item \textbf{Equazioni Differenziali Ordinarie (ODE)}: problemi di Cauchy del I ordine (a variabili separabili con determinazione del tempo di blow-up $T_{\text{max}}$ e lineari con fattore integrante) ed equazioni del II ordine a coefficienti costanti in regime omogeneo e forzato con risonanza.
    \item \textbf{Quesiti Concettuali Rapidi e Test di Parte A}: selezione di quesiti sintetici d'esame su insiemi di definizione, immagini, coordinate polari, curve orarie/velocità ed estremi di insiemi numerici.
\end{enumerate}

\vspace{1em}
\hrule
\vspace{1.5em}

% Inclusione dei file tematici degli esercizi
\section{Limiti, Sviluppi di Taylor, Simboli di Landau e Parti Principali}
\label{sec:esami_limiti_taylor}

\begin{esame}[Guida Metodologica: Sviluppi di Taylor e Parti Principali]
Nel programma del Prof. Alberti, il calcolo della \textbf{parte principale} di un infinitesimo o infinito per $x \to x_0$ richiede di determinare una costante reale $c \neq 0$ e un esponente $\alpha > 0$ tali che:
\[
    f(x) = c (x - x_0)^\alpha + \smallo\left((x - x_0)^\alpha\right) \quad (\text{ovvero } f(x) \sim c (x - x_0)^\alpha)
\]
\textbf{Avvertenze fondamentali per la prova d'esame:}
\begin{enumerate}
    \item \textbf{Cancellazione d'ordine inferiore:} Se i termini di ordine più basso si cancellano esattamente (ad esempio i termini di grado 1 o 2), è obbligatorio proseguire lo sviluppo agli ordini successivi fino a trovare il \emph{primo termine a coefficiente non nullo}.
    \item \textbf{Troncamento corretto nelle composizioni:} Se si calcola $f(g(x))$, occorre sviluppare la funzione interna $g(x)$ fino all'ordine $n$ tale che, una volta sostituita nella funzione esterna, non si perdano potenze rilevanti né si calcolino termini superflui oltre $\smallo(x^n)$.
    \item \textbf{Dipendenza da parametri:} Se il coefficiente del termine dominante si annulla per valori critici del parametro (es. $a = \sqrt{3}$), è indispensabile distinguere i casi ed estrarre il termine successivo dello sviluppo.
\end{enumerate}
\end{esame}

\begin{esempio}{Parte Principale con Parametro e Cancellazione (Compitino 2024/25, II Parte)}{esame_taylor_param1}
Sia dato il parametro reale $a > 0$. Consideriamo la funzione:
\[
    f(x) := \ln(x^3 + a) - \ln(x^3 - a) - \frac{2ax^3}{x^6 + 1}
\]
\begin{enumerate}[label=\alph*)]
    \item Determinare la parte principale di $f(x)$ per $x \to +\infty$ al variare del parametro $a > 0$, rispetto al campione $1/x$.
    \item Determinare la parte principale di $f(x)$ per $x \to +\infty$ qualora l'ultimo addendo sia sostituito con $-\frac{2ax^3}{x^6 - 1}$.
\end{enumerate}
\tcblower
\textbf{Soluzione Dettagliata:}

\paragraph{Quesito (a):} Raccogliamo il termine dominante $x^3$ all'interno dei logaritmi per ricondurci a sviluppi di Maclaurin centrati nell'origine:
\begin{align*}
    f(x) &= \left[ \ln(x^3) + \ln\left(1 + \frac{a}{x^3}\right) \right] - \left[ \ln(x^3) + \ln\left(1 - \frac{a}{x^3}\right) \right] - \frac{2ax^3}{x^6\left(1 + x^{-6}\right)} \\
    &= \ln\left(1 + ax^{-3}\right) - \ln\left(1 - ax^{-3}\right) - 2ax^{-3}\left(1 + x^{-6}\right)^{-1}
\end{align*}
Poniamo le variabili infinitesime $t = ax^{-3} \to 0$ ed $s = x^{-6} \to 0$ per $x \to +\infty$.
Ricordiamo gli sviluppi di Maclaurin fondamentali:
\begin{align*}
    \ln(1 + t) &= t - \frac{t^2}{2} + \frac{t^3}{3} - \frac{t^4}{4} + \frac{t^5}{5} + \bigO(t^6) \\
    \ln(1 - t) &= -t - \frac{t^2}{2} - \frac{t^3}{3} - \frac{t^4}{4} - \frac{t^5}{5} + \bigO(t^6) \\
    \implies \ln(1+t) - \ln(1-t) &= 2t + \frac{2}{3}t^3 + \frac{2}{5}t^5 + \bigO(t^7) \\
    (1 + s)^{-1} &= 1 - s + s^2 + \bigO(s^3)
\end{align*}
Sostituendo $t = ax^{-3}$ ed $s = x^{-6}$:
\begin{align*}
    \ln(1 + ax^{-3}) - \ln(1 - ax^{-3}) &= 2ax^{-3} + \frac{2}{3}a^3 x^{-9} + \frac{2}{5}a^5 x^{-15} + \bigO(x^{-21}) \\
    2ax^{-3}(1 + x^{-6})^{-1} &= 2ax^{-3}\left(1 - x^{-6} + x^{-12} + \bigO(x^{-18})\right) \\
    &= 2ax^{-3} - 2ax^{-9} + 2ax^{-15} + \bigO(x^{-21})
\end{align*}
Sottraendo i due contributi, il termine dominante di ordine $x^{-3}$ si cancella esattamente:
\begin{align*}
    f(x) &= \left(2ax^{-3} + \frac{2}{3}a^3 x^{-9} + \bigO(x^{-15})\right) - \left(2ax^{-3} - 2ax^{-9} + \bigO(x^{-15})\right) \\
    &= \left(\frac{2}{3}a^3 + 2a\right)x^{-9} + \bigO(x^{-15}) = 2a\left(\frac{a^2}{3} + 1\right)x^{-9} + \bigO(x^{-15})
\end{align*}
Poiché $a > 0$, la quantità $2a\left(\frac{a^2}{3} + 1\right)$ è strettamente positiva per ogni $a > 0$ e non si annulla mai.
Pertanto, per ogni $a > 0$:
\[
    \operatorname{p.p.}(f(x)) = 2a\left(\frac{a^2}{3} + 1\right)\frac{1}{x^9} \quad \text{per } x \to +\infty \quad (\text{ordine di infinitesimo } 9)
\]

\paragraph{Quesito (b):} Sostituendo con l'addendo $-\frac{2ax^3}{x^6 - 1} = -2ax^{-3}(1 - x^{-6})^{-1}$, utilizziamo lo sviluppo $(1 - s)^{-1} = 1 + s + s^2 + \bigO(s^3)$:
\begin{align*}
    2ax^{-3}(1 - x^{-6})^{-1} &= 2ax^{-3}\left(1 + x^{-6} + x^{-12} + \bigO(x^{-18})\right) \\
    &= 2ax^{-3} + 2ax^{-9} + 2ax^{-15} + \bigO(x^{-21})
\end{align*}
L'espansione della funzione modificata $\tilde{f}(x)$ diventa:
\begin{align*}
    \tilde{f}(x) &= \left[2ax^{-3} + \frac{2}{3}a^3 x^{-9} + \frac{2}{5}a^5 x^{-15} + \bigO(x^{-21})\right] \\
    &\quad - \left[2ax^{-3} + 2ax^{-9} + 2ax^{-15} + \bigO(x^{-21})\right] \\
    &= \left(\frac{2}{3}a^3 - 2a\right)x^{-9} + \left(\frac{2}{5}a^5 - 2a\right)x^{-15} + \bigO(x^{-21}) \\
    &= 2a\left(\frac{a^2}{3} - 1\right)x^{-9} + 2a\left(\frac{a^4}{5} - 1\right)x^{-15} + \bigO(x^{-21})
\end{align*}
Analizziamo l'annullamento del coefficiente di ordine $x^{-9}$:
\[
    \frac{a^2}{3} - 1 = 0 \iff a^2 = 3 \iff a = \sqrt{3} \quad (\text{essendo } a > 0)
\]
\begin{itemize}
    \item \textbf{Se $a \neq \sqrt{3}$:} il coefficiente di $x^{-9}$ è non nullo, quindi:
    \[
        \operatorname{p.p.}(\tilde{f}(x)) = 2a\left(\frac{a^2}{3} - 1\right)\frac{1}{x^9} \quad (\text{ordine } 9)
    \]
    \item \textbf{Se $a = \sqrt{3}$:} il coefficiente di $x^{-9}$ si annulla identicamente. Il termine dominante è quello di ordine $x^{-15}$, il cui coefficiente vale:
    \[
        2\sqrt{3}\left(\frac{(\sqrt{3})^4}{5} - 1\right) = 2\sqrt{3}\left(\frac{9}{5} - 1\right) = 2\sqrt{3}\cdot\frac{4}{5} = \frac{8\sqrt{3}}{5} \neq 0
    \]
    Quindi per $a = \sqrt{3}$:
    \[
        \operatorname{p.p.}(\tilde{f}(x)) = \frac{8\sqrt{3}}{5}\frac{1}{x^{15}} \quad (\text{ordine di infinitesimo } 15)
    \]
\end{itemize}
\end{esempio}

\begin{esempio}{Polinomi di Taylor di Funzioni Composte (Appello 2023/24)}{esame_taylor_comp}
Calcolare il polinomio di Taylor di ordine 4 centrato in $x_0 = 0$ della funzione:
\[
    g(x) := \cos(\ln(1 + 2x)) - e^{-2x^2}
\]
\tcblower
\textbf{Soluzione Dettagliata:}

\paragraph{1. Sviluppo dell'Argomento Interno:}
Poniamo $u(x) = \ln(1 + 2x)$. Poiché cerchiamo lo sviluppo globale al quarto ordine e $u(0) = 0$ con $u(x) = \bigO(x)$, sviluppiamo il logaritmo fino all'ordine 4:
\[
    u(x) = (2x) - \frac{(2x)^2}{2} + \frac{(2x)^3}{3} - \frac{(2x)^4}{4} + \smallo(x^4) = 2x - 2x^2 + \frac{8}{3}x^3 - 4x^4 + \smallo(x^4)
\]

\paragraph{2. Sviluppo del Coseno:}
Lo sviluppo notevole di $\cos(u)$ per $u \to 0$ è:
\[
    \cos(u) = 1 - \frac{u^2}{2!} + \frac{u^4}{4!} + \smallo(u^4) = 1 - \frac{u^2}{2} + \frac{u^4}{24} + \smallo(u^4)
\]
Calcoliamo le potenze di $u(x)$ trascurando tutti i termini di grado superiore a 4:
\begin{align*}
    u^2(x) &= \left(2x - 2x^2 + \frac{8}{3}x^3 + \smallo(x^3)\right)^2 \\
    &= (2x)^2 + (-2x^2)^2 + 2(2x)(-2x^2) + 2(2x)\left(\frac{8}{3}x^3\right) + \smallo(x^4) \\
    &= 4x^2 + 4x^4 - 8x^3 + \frac{32}{3}x^4 + \smallo(x^4) = 4x^2 - 8x^3 + \frac{44}{3}x^4 + \smallo(x^4)
\end{align*}
Per la potenza quarta, poiché $u(x) = 2x + \bigO(x^2)$:
\[
    u^4(x) = (2x)^4 + \smallo(x^4) = 16x^4 + \smallo(x^4)
\]
Sostituendo nello sviluppo del coseno:
\begin{align*}
    \cos(\ln(1 + 2x)) &= 1 - \frac{1}{2}\left(4x^2 - 8x^3 + \frac{44}{3}x^4\right) + \frac{1}{24}(16x^4) + \smallo(x^4) \\
    &= 1 - 2x^2 + 4x^3 - \frac{22}{3}x^4 + \frac{2}{3}x^4 + \smallo(x^4) \\
    &= 1 - 2x^2 + 4x^3 - \frac{20}{3}x^4 + \smallo(x^4)
\end{align*}

\paragraph{3. Sviluppo dell'Esponenziale:}
Posto $v = -2x^2 \to 0$, si ha:
\[
    e^{-2x^2} = 1 + (-2x^2) + \frac{(-2x^2)^2}{2!} + \smallo(x^4) = 1 - 2x^2 + 2x^4 + \smallo(x^4)
\]

\paragraph{4. Combinazione Lineare e Polinomio di Taylor:}
Sottraendo i due sviluppi asintotici:
\begin{align*}
    g(x) &= \left(1 - 2x^2 + 4x^3 - \frac{20}{3}x^4 + \smallo(x^4)\right) - \left(1 - 2x^2 + 2x^4 + \smallo(x^4)\right) \\
    &= (1 - 1) + (-2 - (-2))x^2 + 4x^3 + \left(-\frac{20}{3} - 2\right)x^4 + \smallo(x^4) \\
    &= 4x^3 - \frac{26}{3}x^4 + \smallo(x^4)
\end{align*}
Il polinomio di Maclaurin di grado 4 è pertanto:
\[
    P_4(x) = 4x^3 - \frac{26}{3}x^4
\]
\end{esempio}

\begin{esempio}{Gerarchia Asintotica e Relazione $\ll$ (Compitino 2024/25)}{esame_landau_ord}
Ordinare le seguenti funzioni rispetto alla relazione di trascurabilità $\ll$ per $x \to +\infty$:
\[
    f_1(x) = \frac{x^2 + 1}{x^2 + 3\ln x}, \qquad f_2(x) = \frac{x^{-3} + x^3}{e^{3x} + 1}, \qquad f_3(x) = \frac{1 - 2x^2}{x\ln x}, \qquad f_4(x) = \frac{e^{2x} - 1}{x^4}
\]
\tcblower
\textbf{Soluzione Dettagliata:}
Ricordiamo che per definizione $f(x) \ll g(x)$ per $x \to +\infty$ se e solo se $\lim_{x \to +\infty} \frac{f(x)}{g(x)} = 0$.
Determiniamo il comportamento asintotico dominante di ciascuna funzione:
\begin{enumerate}
    \item $f_1(x) = \frac{x^2(1 + x^{-2})}{x^2(1 + 3x^{-2}\ln x)} \xrightarrow{x \to +\infty} 1 \implies f_1(x) = \Theta(1)$ (ordine costante finito e non nullo).
    \item $f_2(x) = \frac{x^3(1 + x^{-6})}{e^{3x}(1 + e^{-3x})} \sim \frac{x^3}{e^{3x}} = x^3 e^{-3x} \xrightarrow{x \to +\infty} 0$ (infinitesimo di tipo esponenziale).
    \item $f_3(x) = \frac{-2x^2(1 - \frac{1}{2x^2})}{x\ln x} \sim -\frac{2x}{\ln x} \xrightarrow{x \to +\infty} -\infty$, con modulo $|f_3(x)| \sim \frac{2x}{\ln x}$ (infinito sub-lineare).
    \item $f_4(x) = \frac{e^{2x}(1 - e^{-2x})}{x^4} \sim \frac{e^{2x}}{x^4} \xrightarrow{x \to +\infty} +\infty$ (infinito di tipo esponenziale).
\end{enumerate}
Confrontiamo i rapporti a due a due:
\begin{itemize}
    \item $\lim_{x \to +\infty} \frac{f_2(x)}{f_1(x)} = \lim_{x \to +\infty} \frac{x^3 e^{-3x}}{1} = 0 \implies f_2(x) \ll f_1(x)$.
    \item $\lim_{x \to +\infty} \frac{f_1(x)}{f_3(x)} = \lim_{x \to +\infty} \frac{1}{-2x/\ln x} = \lim_{x \to +\infty} \left(-\frac{\ln x}{2x}\right) = 0 \implies f_1(x) \ll f_3(x)$.
    \item $\lim_{x \to +\infty} \frac{f_3(x)}{f_4(x)} = \lim_{x \to +\infty} \frac{-2x/\ln x}{e^{2x}/x^4} = \lim_{x \to +\infty} \left(-\frac{2x^5}{e^{2x}\ln x}\right) = 0 \implies f_3(x) \ll f_4(x)$.
\end{itemize}
Ne deduciamo la catena ordinata completa:
\[
    f_2(x) \ll f_1(x) \ll f_3(x) \ll f_4(x) \quad \text{per } x \to +\infty
\]
\end{esempio}

\begin{esempio}{Forma Indeterminata con Parametro Reale (Appello 2023/24)}{esame_limite_param_forma_indet}
Calcolare al variare del parametro $\alpha > 0$ il limite:
\[
    L(\alpha) := \lim_{x \to 0^+} \frac{e^{x^2} - \cos(\sqrt{2}x) - 2x^2}{\ln(1 + x^4) + x^\alpha}
\]
\tcblower
\textbf{Soluzione Dettagliata:}

\paragraph{1. Sviluppo del Numeratore $N(x)$:}
Per $x \to 0$, sviluppiamo $e^{x^2}$ e $\cos(\sqrt{2}x)$ fino all'ordine 4:
\begin{align*}
    e^{x^2} &= 1 + x^2 + \frac{(x^2)^2}{2!} + \bigO(x^6) = 1 + x^2 + \frac{x^4}{2} + \bigO(x^6) \\
    \cos(\sqrt{2}x) &= 1 - \frac{(\sqrt{2}x)^2}{2!} + \frac{(\sqrt{2}x)^4}{4!} + \bigO(x^6) \\
    &= 1 - x^2 + \frac{4x^4}{24} + \bigO(x^6) = 1 - x^2 + \frac{x^4}{6} + \bigO(x^6)
\end{align*}
Sostituendo nel numeratore $N(x)$:
\begin{align*}
    N(x) &= \left(1 + x^2 + \frac{x^4}{2} + \bigO(x^6)\right) - \left(1 - x^2 + \frac{x^4}{6} + \bigO(x^6)\right) - 2x^2 \\
    &= (1 - 1) + (x^2 + x^2 - 2x^2) + \left(\frac{1}{2} - \frac{1}{6}\right)x^4 + \bigO(x^6) \\
    &= \frac{3 - 1}{6}x^4 + \bigO(x^6) = \frac{1}{3}x^4 + \bigO(x^6)
\end{align*}
Dunque $N(x) \sim \frac{1}{3}x^4$ per $x \to 0^+$.

\paragraph{2. Sviluppo del Denominatore $D(x)$ al variare di $\alpha > 0$:}
Poiché $\ln(1 + x^4) = x^4 + \bigO(x^8)$, il denominatore si scrive:
\[
    D(x) = x^4 + \bigO(x^8) + x^\alpha
\]
Il comportamento asintotico di $D(x)$ per $x \to 0^+$ è determinato dal confronto tra le potenze $x^4$ e $x^\alpha$:
\begin{itemize}
    \item \textbf{Caso $0 < \alpha < 4$:} la potenza $x^\alpha$ è di grado inferiore rispetto a $x^4$, quindi $x^\alpha$ è dominante ($x^4 \ll x^\alpha$).
    Dunque $D(x) \sim x^\alpha$, da cui:
    \[
        \frac{N(x)}{D(x)} \sim \frac{\frac{1}{3}x^4}{x^\alpha} = \frac{1}{3}x^{4 - \alpha} \xrightarrow{x \to 0^+} 0 \quad (\text{essendo } 4 - \alpha > 0)
    \]
    \item \textbf{Caso $\alpha = 4$:} i due termini hanno lo stesso ordine di infinitesimo:
    \[
        D(x) = x^4 + x^4 + \bigO(x^8) = 2x^4 + \bigO(x^8) \sim 2x^4
    \]
    Pertanto:
    \[
        L(4) = \lim_{x \to 0^+} \frac{\frac{1}{3}x^4 + \bigO(x^6)}{2x^4 + \bigO(x^8)} = \frac{1/3}{2} = \frac{1}{6}
    \]
    \item \textbf{Caso $\alpha > 4$:} la potenza $x^4$ è dominante rispetto a $x^\alpha$ ($x^\alpha \ll x^4$).
    Dunque $D(x) \sim x^4$, da cui:
    \[
        L(\alpha) = \lim_{x \to 0^+} \frac{\frac{1}{3}x^4 + \bigO(x^6)}{x^4 + \bigO(x^\alpha)} = \frac{1}{3}
    \]
\end{itemize}
\paragraph{Sintesi del Risultato:}
\[
    L(\alpha) = \begin{cases}
        0 & \text{se } 0 < \alpha < 4 \\
        \frac{1}{6} & \text{se } \alpha = 4 \\
        \frac{1}{3} & \text{se } \alpha > 4
    \end{cases}
\]
\end{esempio}

\section{Calcolo Differenziale, Punti Critici e Studio di Funzione}
\label{sec:esami_studio_funzione}

\begin{esame}[Protocollo Rigoroso per lo Studio Qualitativo e Punti Singolari]
Nelle prove d'esame del Prof. Alberti, lo studio qualitativo di funzione richiede massima accuratezza nella gestione dei punti singolari e del calcolo differenziale:
\begin{enumerate}
    \item \textbf{Classificazione delle Singolarità della Derivata Prima:}
    Sia $f$ continua in $x_0$. Calcolando i limiti destro e sinistro di $f'(x)$ per $x \to x_0^\pm$:
    \begin{itemize}
        \item \textbf{Punto Angoloso:} $f'(x_0^+) \neq f'(x_0^-)$ con almeno uno dei due limiti finito.
        \item \textbf{Cuspide:} $f'(x_0^+) = \pm \infty$ e $f'(x_0^-) = \mp \infty$ (segni discordi infiniti).
        \item \textbf{Flesso a Tangente Verticale:} $f'(x_0^+) = f'(x_0^-) = +\infty$ oppure $-\infty$ (segni concordi infiniti).
    \end{itemize}
    \item \textbf{Posizione del Grafico rispetto agli Asintoti:}
    Determinata l'equazione dell'asintoto obliquo $y = mx + q$, studiare sempre il segno dello scarto asintotico $f(x) - (mx + q)$ al primo ordine non nullo per $x \to \pm \infty$.
    \item \textbf{Ottimizzazione su Compatti:}
    Per il Teorema di Weierstrass, il massimo e minimo assoluto su un intervallo compatto $[a, b]$ vanno ricercati confrontando: i punti critici interni ($f'(x) = 0$), i punti di non derivabilità interni e i due estremi di bordo $\{a, b\}$.
\end{enumerate}
\end{esame}

\begin{esempio}{Studio Completo di Funzione con Radici, Asintoti e Flessi (Appello 2023/24)}{esame_studio_f1}
Studiare qualitativamente la funzione:
\[
    f(x) := \sqrt[3]{x^3 - 3x^2}
\]
determinandone dominio naturale, limiti agli estremi, eventuali asintoti (verticali, orizzontali, obliqui), derivabilità, classificazione dei punti di non derivabilità, intervalli di monotonia, estremi locali ed assoluti, derivata seconda e concavità/convessità.
\tcblower
\textbf{Soluzione Dettagliata:}

\paragraph{1. Dominio, Zeri e Segno:}
L'indice della radice è $3$ (dispari), pertanto la funzione è definita e continua su tutto l'asse reale:
\[
    \dom(f) = \R = (-\infty, +\infty)
\]
Intersezioni con gli assi e segno:
\begin{itemize}
    \item $f(x) = 0 \iff x^3 - 3x^2 = 0 \iff x^2(x - 3) = 0 \implies x = 0$ (radice doppia) e $x = 3$ (radice semplice).
    \item Poiché $x^2 \ge 0$, il segno coincide con quello di $x - 3$:
    \[
        f(x) > 0 \iff x \in (3, +\infty), \qquad f(x) < 0 \iff x \in (-\infty, 0) \cup (0, 3)
    \]
\end{itemize}

\paragraph{2. Limiti agli Estremi e Asintoti Obliqui:}
Per $x \to \pm \infty$:
\[
    \lim_{x \to \pm \infty} f(x) = \lim_{x \to \pm \infty} x \left(1 - \frac{3}{x}\right)^{1/3} = \pm \infty
\]
Ricerca di asintoti obliqui $y = mx + q$:
\begin{align*}
    m &= \lim_{x \to \pm \infty} \frac{f(x)}{x} = \lim_{x \to \pm \infty} \left(1 - \frac{3}{x}\right)^{1/3} = 1 \\
    q &= \lim_{x \to \pm \infty} [f(x) - x] = \lim_{x \to \pm \infty} x\left[\left(1 - \frac{3}{x}\right)^{1/3} - 1\right]
\end{align*}
Sviluppiamo con la formula binomiale $(1 + t)^{1/3} = 1 + \frac{1}{3}t - \frac{1}{9}t^2 + \bigO(t^3)$ con $t = -3/x \to 0$:
\[
    \left(1 - \frac{3}{x}\right)^{1/3} = 1 + \frac{1}{3}\left(-\frac{3}{x}\right) - \frac{1}{9}\left(-\frac{3}{x}\right)^2 + \bigO(x^{-3}) = 1 - \frac{1}{x} - \frac{1}{x^2} + \bigO(x^{-3})
\]
Dunque:
\[
    q = \lim_{x \to \pm \infty} x\left(-\frac{1}{x} - \frac{1}{x^2} + \bigO(x^{-3})\right) = \lim_{x \to \pm \infty} \left(-1 - \frac{1}{x} + \bigO(x^{-2})\right) = -1
\]
La retta $y = x - 1$ è \textbf{asintoto obliquo completo} sia a $+\infty$ che a $-\infty$.
Inoltre, lo scarto $f(x) - (x - 1) = -\frac{1}{x} + \bigO(x^{-2})$ indica la posizione relativa:
\begin{itemize}
    \item Per $x \to +\infty$: $f(x) - (x - 1) \sim -\frac{1}{x} < 0 \implies$ grafico al di \textbf{sotto} dell'asintoto.
    \item Per $x \to -\infty$: $f(x) - (x - 1) \sim -\frac{1}{x} > 0 \implies$ grafico al di \textbf{sopra} dell'asintoto.
\end{itemize}

\paragraph{3. Derivata Prima e Punti di Non Derivabilità:}
Scrivendo $f(x) = (x^3 - 3x^2)^{1/3}$, per $x \notin \{0, 3\}$ applichiamo la regola della catena:
\[
    f'(x) = \frac{1}{3}(x^3 - 3x^2)^{-2/3} (3x^2 - 6x) = \frac{3x(x - 2)}{3 (x^2(x - 3))^{2/3}} = \frac{x(x - 2)}{x^{4/3}(x - 3)^{2/3}} = \frac{x - 2}{x^{1/3}(x - 3)^{2/3}}
\]
Studiamo i limiti di $f'(x)$ nei punti di annullamento dell'argomento della radice:
\begin{itemize}
    \item In $x = 0$: per $x \to 0$, il fattore $(x-3)^{2/3} \to (-3)^{2/3} = 3^{2/3} > 0$ e $x - 2 \to -2 < 0$. Dunque:
    \[
        \lim_{x \to 0^+} f'(x) = \frac{-2}{0^+ \cdot 3^{2/3}} = -\infty, \qquad \lim_{x \to 0^-} f'(x) = \frac{-2}{0^- \cdot 3^{2/3}} = +\infty
    \]
    Il punto $(0, 0)$ è un punto di \textbf{cuspide} rivolta verso l'alto.
    \item In $x = 3$: per $x \to 3$, $x - 2 \to 1 > 0$ e $x^{1/3} \to 3^{1/3} > 0$, mentre $(x-3)^{2/3} \to 0^+$ sia da destra che da sinistra:
    \[
        \lim_{x \to 3^+} f'(x) = \lim_{x \to 3^-} f'(x) = \frac{1}{3^{1/3} \cdot 0^+} = +\infty
    \]
    Il punto $(3, 0)$ è un punto di \textbf{flesso a tangente verticale ascendente}.
\end{itemize}

\paragraph{4. Monotonia ed Estremi Relativi:}
Poiché $(x-3)^{2/3} > 0$ per ogni $x \neq 3$, il segno di $f'(x)$ coincide con il segno di $\frac{x-2}{x^{1/3}}$:
\begin{itemize}
    \item Su $(-\infty, 0)$: $x - 2 < 0$ e $x^{1/3} < 0 \implies f'(x) > 0$ ($f$ strettamente crescente).
    \item Su $(0, 2)$: $x - 2 < 0$ e $x^{1/3} > 0 \implies f'(x) < 0$ ($f$ strettamente decrescente).
    \item In $x = 2$: $f'(2) = 0$ (punto stazionario).
    \item Su $(2, 3) \cup (3, +\infty)$: $x - 2 > 0$ e $x^{1/3} > 0 \implies f'(x) > 0$ ($f$ strettamente crescente).
\end{itemize}
Classificazione degli estremi:
\begin{itemize}
    \item $x = 0$ è punto di \textbf{massimo locale} con valore $f(0) = 0$.
    \item $x = 2$ è punto di \textbf{minimo locale} con valore $f(2) = \sqrt[3]{8 - 12} = -\sqrt[3]{4} \approx -1.587$.
    \item Poiché $\lim_{x \to \pm \infty} f(x) = \pm \infty$, $\Imm(f) = \R$ e non esistono massimo né minimo assoluti su $\R$.
\end{itemize}

\paragraph{5. Derivata Seconda e Concavità/Convessità:}
Calcoliamo $f''(x)$ per $x \notin \{0, 3\}$ derivando $f'(x) = (x-2) x^{-1/3} (x-3)^{-2/3}$:
Usando la derivata logaritmica $\frac{f''(x)}{f'(x)} = \frac{1}{x-2} - \frac{1}{3x} - \frac{2}{3(x-3)}$:
\begin{align*}
    \frac{f''(x)}{f'(x)} &= \frac{1}{x-2} - \frac{x-3 + 2x}{3x(x-3)} = \frac{1}{x-2} - \frac{x-1}{x(x-3)} \\
    &= \frac{x(x-3) - (x-2)(x-1)}{x(x-2)(x-3)} = \frac{-2}{x(x-2)(x-3)}
\end{align*}
Moltiplicando per $f'(x) = \frac{x-2}{x^{1/3}(x-3)^{2/3}}$:
\[
    f''(x) = \frac{x-2}{x^{1/3}(x-3)^{2/3}} \cdot \frac{-2}{x(x-2)(x-3)} = \frac{-2}{x^{4/3}(x-3)^{5/3}}
\]
Analisi del segno di $f''(x)$:
\begin{itemize}
    \item Il termine al numeratore è la costante $-2 < 0$.
    \item Il fattore $x^{4/3} = (x^2)^{2/3} > 0$ per ogni $x \neq 0$.
    \item Il fattore $(x-3)^{5/3}$ ha lo stesso segno di $x-3$:
    \begin{itemize}
        \item Per $x < 3$ ($x \neq 0$): $x - 3 < 0 \implies (x-3)^{5/3} < 0 \implies f''(x) > 0$.
        La funzione è \textbf{strettamente convessa} su $(-\infty, 0)$ e su $(0, 3)$.
        \item Per $x > 3$: $x - 3 > 0 \implies (x-3)^{5/3} > 0 \implies f''(x) < 0$.
        La funzione è \textbf{strettamente concava} su $(3, +\infty)$.
    \end{itemize}
\end{itemize}
In $x = 3$ la funzione cambia concavità (da convessa a concava), confermando rigorosamente la presenza del \textbf{punto di flesso a tangente verticale} $(3, 0)$.
\end{esempio}

\begin{esempio}{Massimi e Minimi con Parametro su Insieme Compatto (Appello 2022/23)}{esame_max_min_param}
Data la funzione $f(x) := x^2 e^{-ax}$ con parametro reale $a > 0$, determinare:
\begin{enumerate}[label=\alph*)]
    \item I punti di massimo e minimo assoluti di $f$ sull'intervallo compatto $[0, 2]$ al variare di $a > 0$.
    \item Il valore massimo $M(a) = \max_{x \in [0, 2]} f(x)$ e la regolarità di $M(a)$ rispetto ad $a$.
\end{enumerate}
\tcblower
\textbf{Soluzione Dettagliata:}
La funzione $f$ è di classe $C^\infty(\R)$. Sull'intervallo compatto $[0, 2]$, per il Teorema di Weierstrass, $f$ ammette sicuramente massimo e minimo assoluti.

\paragraph{1. Punti Stazionari ed Estremi di Bordo:}
Calcoliamo la derivata prima:
\[
    f'(x) = 2x e^{-ax} + x^2(-a e^{-ax}) = x(2 - ax)e^{-ax}
\]
I punti stazionari in $\R$ sono:
\[
    x_1 = 0, \qquad x_2 = \frac{2}{a}
\]
Valori nei punti notevoli:
\begin{itemize}
    \item In $x = 0$: $f(0) = 0$. Poiché $f(x) = x^2 e^{-ax} \ge 0$ per ogni $x \ge 0$, il punto $x = 0$ è sempre il punto di \textbf{minimo assoluto}, con $\min_{[0, 2]} f = 0$.
    \item In $x = 2$: $f(2) = 4 e^{-2a}$.
    \item In $x = \frac{2}{a}$: $f\left(\frac{2}{a}\right) = \left(\frac{2}{a}\right)^2 e^{-a(2/a)} = \frac{4}{a^2 e^2}$.
\end{itemize}

\paragraph{2. Posizione del Punto Critico rispetto all'Intervallo $[0, 2]$:}
\begin{itemize}
    \item \textbf{Caso 1: $\frac{2}{a} \ge 2 \iff 0 < a \le 1$:}
    Il punto stazionario interno non appartiene a $(0, 2)$. Per ogni $x \in (0, 2)$ si ha $2 - ax > 0$, dunque $f'(x) > 0$.
    La funzione è strettamente crescente su $[0, 2]$.
    Il punto di massimo assoluto è l'estremo destro:
    \[
        x_{\max} = 2, \qquad M(a) = f(2) = 4e^{-2a}
    \]
    \item \textbf{Caso 2: $\frac{2}{a} < 2 \iff a > 1$:}
    Il punto $x_c = \frac{2}{a}$ cade all'interno dell'intervallo $(0, 2)$.
    Essendo $f'(x) > 0$ per $0 < x < \frac{2}{a}$ e $f'(x) < 0$ per $\frac{2}{a} < x < 2$, il punto $x_c$ è il punto di massimo assoluto su $[0, 2]$:
    \[
        x_{\max} = \frac{2}{a}, \qquad M(a) = f\left(\frac{2}{a}\right) = \frac{4}{a^2 e^2}
    \]
\end{itemize}

\paragraph{3. Studio della Funzione Valore Massimo $M(a)$:}
La funzione massimo $M \colon (0, +\infty) \to \R$ è data da:
\[
    M(a) = \begin{cases}
        4e^{-2a} & \text{se } 0 < a \le 1 \\
        \frac{4}{e^2 a^2} & \text{se } a > 1
    \end{cases}
\]
Verifichiamo la continuità e derivabilità nel punto di giunzione $a = 1$:
\begin{itemize}
    \item Continuità: $M(1^-) = 4e^{-2}$, $M(1^+) = \frac{4}{e^2 \cdot 1^2} = 4e^{-2} = M(1)$. La funzione è continua.
    \item Derivabilità:
    \[
        M'(a) = \begin{cases}
            -8e^{-2a} & \text{per } 0 < a < 1 \implies M'(1^-) = -8e^{-2} \\
            -\frac{8}{e^2 a^3} & \text{per } a > 1 \implies M'(1^+) = -\frac{8}{e^2 \cdot 1^3} = -8e^{-2}
        \end{cases}
    \]
    Poiché le derivate destra e sinistra coincidono, $M(a)$ è di classe $C^1((0, +\infty))$.
\end{itemize}
\end{esempio}

\begin{esempio}{Equazioni Parametriche e Teorema dei Valori Intermedi (Appello 2024/25)}{esame_studio_soluzioni_param}
Determinare al variare del parametro reale $\lambda \in \R$ il numero di soluzioni reali distinte dell'equazione:
\[
    (x^2 - 4)e^{-x} = \lambda
\]
\tcblower
\textbf{Soluzione Dettagliata:}
Poniamo $h(x) := (x^2 - 4)e^{-x}$. Risolvere l'equazione equivale a intersecare il grafico di $h(x)$ con la retta orizzontale $y = \lambda$.

\paragraph{1. Limiti agli Estremi e Zeri:}
\begin{itemize}
    \item $\dom(h) = \R$. Intersezioni asse $x$: $h(x) = 0 \iff x = \pm 2$.
    \item $\lim_{x \to +\infty} (x^2 - 4)e^{-x} = 0^+$ (per la gerarchia infinitesima dell'esponenziale).
    \item $\lim_{x \to -\infty} (x^2 - 4)e^{-x} = (+\infty)(+\infty) = +\infty$.
\end{itemize}
La retta $y = 0$ è asintoto orizzontale destro per $x \to +\infty$.

\paragraph{2. Punti Stazionari e Monotonia:}
Calcoliamo la derivata prima:
\[
    h'(x) = 2x e^{-x} - (x^2 - 4)e^{-x} = -(x^2 - 2x - 4)e^{-x}
\]
Le radici del trinomio $x^2 - 2x - 4 = 0$ sono $x_{1,2} = 1 \pm \sqrt{5}$:
\[
    x_1 = 1 - \sqrt{5} \approx -1.236, \qquad x_2 = 1 + \sqrt{5} \approx 3.236
\]
Studio del segno di $h'(x)$ (tenendo conto del segno meno davanti al trinomio):
\begin{itemize}
    \item Per $x < 1 - \sqrt{5}$: $h'(x) < 0 \implies h$ è strettamente decrescente da $+\infty$ a $h(x_1)$.
    \item Per $1 - \sqrt{5} < x < 1 + \sqrt{5}$: $h'(x) > 0 \implies h$ è strettamente crescente da $h(x_1)$ a $h(x_2)$.
    \item Per $x > 1 + \sqrt{5}$: $h'(x) < 0 \implies h$ è strettamente decrescente da $h(x_2)$ a $0^+$.
\end{itemize}
I valori estremi locali sono:
\begin{align*}
    h(x_1) &= \left((1-\sqrt{5})^2 - 4\right)e^{-(1-\sqrt{5})} = (2 - 2\sqrt{5})e^{\sqrt{5}-1} < 0 \quad (\text{minimo assoluto}) \\
    h(x_2) &= \left((1+\sqrt{5})^2 - 4\right)e^{-(1+\sqrt{5})} = (2 + 2\sqrt{5})e^{-(\sqrt{5}+1)} > 0 \quad (\text{massimo locale})
\end{align*}

\paragraph{3. Discussione del Numero di Soluzioni al variare di $\lambda$:}
Applicando il Teorema dei Valori Intermedi e la stretta monotonia su ciascuno dei tre intervalli $(-\infty, x_1]$, $[x_1, x_2]$ e $[x_2, +\infty)$:
\begin{itemize}
    \item Per $\lambda < h(x_1)$: \textbf{0 soluzioni} reali.
    \item Per $\lambda = h(x_1)$: \textbf{1 soluzione} reale ($x = x_1$).
    \item Per $h(x_1) < \lambda \le 0$: \textbf{2 soluzioni} reali distinte (una in $(-\infty, x_1)$ e una in $(x_1, 2]$).
    \item Per $0 < \lambda < h(x_2)$: \textbf{3 soluzioni} reali distinte (una in $(-\infty, x_1)$, una in $(2, x_2)$ e una in $(x_2, +\infty)$).
    \item Per $\lambda = h(x_2)$: \textbf{2 soluzioni} reali distinte (una in $(-\infty, x_1)$ e una doppia stazionaria in $x = x_2$).
    \item Per $\lambda > h(x_2)$: \textbf{1 soluzione} reale (in $(-\infty, x_1)$).
\end{itemize}
\end{esempio}

\section{Calcolo Integrale, Primitive, Aree, Volumi e Integrali Impropri}
\label{sec:esami_integrali}

\begin{esame}[Strategie d'Esame per Integrali Impropri, Aree e Volumi]
Nelle prove scritte del Prof. Alberti sul calcolo integrale:
\begin{enumerate}
    \item \textbf{Criteri per Integrali Impropri Parametrici:}
    Separare sempre le singolarità isolate (es. $x \to 0^+$, $x \to +\infty$, punti interni di discontinuità) e studiare separatamente ciascun integrale parziale.
    Ricordare gli integrali campione:
    \[
        \int_0^1 \frac{1}{x^p |\ln x|^q} \dx < \infty \iff (p < 1) \lor (p = 1 \land q > 1)
    \]
    \[
        \int_1^{+\infty} \frac{1}{x^p (\ln x)^q} \dx < \infty \iff (p > 1) \lor (p = 1 \land q > 1)
    \]
    \item \textbf{Teorema di Pappo-Guldino e Condizioni di Applicabilità:}
    Il volume generato dalla rotazione completa di $2\pi$ di una figura piana misurabile $A$ attorno a un asse $r$ complanare è dato da $\operatorname{Vol} = 2\pi d \operatorname{Area}(A)$ (con $d = \operatorname{dist}(G, r)$) \textbf{a patto che l'asse $r$ non intersechi l'interno di $A$}. Se l'asse è secante, la formula non si applica globalmente a causa della sovrapposizione dei volumi e occorre spezzare l'integrale.
\end{enumerate}
\end{esame}

\begin{esempio}{Aree Piane e Volumi di Solidi di Rotazione (Compitino 2024/25, II Parte)}{esame_area_volumi}
Sia $A$ la regione piana limitata definita dall'insieme dei punti $(x, y) \in \R^2$ tali che:
\[
    1 \le y \le \frac{4}{1 + x^2}
\]
\begin{enumerate}[label=\alph*)]
    \item Calcolare l'area della regione $A$.
    \item Calcolare il volume del solido $V_x$ ottenuto ruotando la regione $A$ di un giro completo attorno all'asse delle ascisse ($y = 0$).
    \item Calcolare il volume del solido $V_2$ ottenuto ruotando la regione $A$ di un giro completo attorno alla retta verticale esterna $x = 2$.
\end{enumerate}
\tcblower
\textbf{Soluzione Dettagliata:}

\paragraph{Analisi Geometrica e Dominio di Integrazione:}
La funzione $f(x) := \frac{4}{1 + x^2}$ è continua, pari, strettamente positiva e limitata superiormente da $f(0) = 4$.
I punti di intersezione con la retta $y = 1$ si ottengono risolvendo:
\[
    \frac{4}{1 + x^2} = 1 \iff 1 + x^2 = 4 \iff x^2 = 3 \iff x = \pm\sqrt{3}
\]
La regione $A$ è compresa tra $y = 1$ (bordo inferiore) e $y = f(x)$ (bordo superiore) per $x \in [-\sqrt{3}, \sqrt{3}]$.

\paragraph{Quesito (a): Calcolo dell'Area di $A$}
Sfruttando la simmetria (parità) del dominio e dell'integrando:
\begin{align*}
    \operatorname{Area}(A) &= \int_{-\sqrt{3}}^{\sqrt{3}} \left(\frac{4}{1 + x^2} - 1\right) \dx = 2 \int_0^{\sqrt{3}} \left(\frac{4}{1 + x^2} - 1\right) \dx \\
    &= 2 \Big[ 4 \arctg(x) - x \Big]_0^{\sqrt{3}} = 2 \left( 4 \arctg(\sqrt{3}) - \sqrt{3} \right) \\
    &= 2 \left( 4 \cdot \frac{\pi}{3} - \sqrt{3} \right) = \frac{8\pi}{3} - 2\sqrt{3}
\end{align*}

\paragraph{Quesito (b): Volume del Solido di Rotazione attorno all'Asse $x$}
Applicando il metodo delle sezioni circolari con fette perpendicolari all'asse $x$:
\begin{align*}
    \operatorname{Vol}(V_x) &= \pi \int_{-\sqrt{3}}^{\sqrt{3}} \left( [f(x)]^2 - 1^2 \right) \dx = 2\pi \int_0^{\sqrt{3}} \left( \frac{16}{(1 + x^2)^2} - 1 \right) \dx \\
    &= 32\pi \int_0^{\sqrt{3}} \frac{1}{(1 + x^2)^2} \dx - 2\pi\sqrt{3}
\end{align*}
Calcoliamo la primitiva $\int \frac{1}{(1 + x^2)^2}\dx$ integrando per parti:
\begin{align*}
    \int \frac{1}{1 + x^2}\dx &= \frac{x}{1 + x^2} - \int x \left(-\frac{2x}{(1 + x^2)^2}\right)\dx = \frac{x}{1 + x^2} + 2\int \frac{(x^2 + 1) - 1}{(1 + x^2)^2}\dx \\
    &= \frac{x}{1 + x^2} + 2\int \frac{1}{1 + x^2}\dx - 2\int \frac{1}{(1 + x^2)^2}\dx
\end{align*}
Risolvendo rispetto all'integrale incognito:
\[
    \int \frac{1}{(1 + x^2)^2}\dx = \frac{1}{2}\left( \frac{x}{1 + x^2} + \arctg(x) \right)
\]
Valutando tra $0$ e $\sqrt{3}$:
\[
    \int_0^{\sqrt{3}} \frac{1}{(1 + x^2)^2}\dx = \frac{1}{2}\left( \frac{\sqrt{3}}{1 + 3} + \arctg(\sqrt{3}) \right) = \frac{1}{2}\left( \frac{\sqrt{3}}{4} + \frac{\pi}{3} \right) = \frac{\sqrt{3}}{8} + \frac{\pi}{6}
\]
Sostituendo nell'espressione del volume:
\[
    \operatorname{Vol}(V_x) = 32\pi \left(\frac{\pi}{6} + \frac{\sqrt{3}}{8}\right) - 2\pi\sqrt{3} = \frac{16\pi^2}{3} + 4\pi\sqrt{3} - 2\pi\sqrt{3} = \frac{16\pi^2}{3} + 2\pi\sqrt{3}
\]

\paragraph{Quesito (c): Volume del Solido di Rotazione attorno alla Retta $x = 2$}
Poiché la regione $A$ è contenuta nell'intervallo $[-\sqrt{3}, \sqrt{3}]$ e $\sqrt{3} \approx 1.732 < 2$, l'asse di rotazione $x = 2$ è \textbf{strettamente esterno} alla regione $A$.
Possiamo quindi applicare direttamente il \textbf{Primo Teorema di Pappo-Guldino}:
\begin{itemize}
    \item Essendo $A$ simmetrica rispetto all'asse delle ordinate ($x = 0$), l'ascissa del baricentro è $x_G = 0$.
    \item La distanza del baricentro dall'asse di rotazione $x = 2$ è $d = |x_G - 2| = 2$.
\end{itemize}
Il volume del solido generato da una rotazione di $2\pi$ è:
\[
    \operatorname{Vol}(V_2) = 2\pi \cdot d \cdot \operatorname{Area}(A) = 2\pi \cdot 2 \cdot \left(\frac{8\pi}{3} - 2\sqrt{3}\right) = \frac{32\pi^2}{3} - 8\pi\sqrt{3}
\]
\textit{Verifica analitica mediante gusci cilindrici:}
\begin{align*}
    \operatorname{Vol}(V_2) &= 2\pi \int_{-\sqrt{3}}^{\sqrt{3}} (2 - x)(f(x) - 1)\dx \\
    &= 4\pi \int_{-\sqrt{3}}^{\sqrt{3}} (f(x) - 1)\dx - 2\pi \int_{-\sqrt{3}}^{\sqrt{3}} x(f(x) - 1)\dx \\
    &= 4\pi \cdot \operatorname{Area}(A) - 0 = \frac{32\pi^2}{3} - 8\pi\sqrt{3}
\end{align*}
\end{esempio}

\begin{esempio}{Integrale Improprio con Parametro Reale (Appello 2023/24)}{esame_integrale_improprio_param}
Discutere al variare del parametro reale $\alpha \in \R$ la convergenza dell'integrale improprio:
\[
    I(\alpha) := \int_0^{+\infty} \frac{\ln(1 + x^\alpha)}{x^2 \sqrt{1 + x}} \dx
\]
\tcblower
\textbf{Soluzione Dettagliata:}
L'integrando $f_\alpha(x) = \frac{\ln(1 + x^\alpha)}{x^2 \sqrt{1 + x}}$ è continuo e a valori strettamente positivi su $(0, +\infty)$.
Le singolarità sono isolate in $x = 0^+$ e per $x \to +\infty$. Spezziamo l'integrale nel punto neutro $x = 1$:
\[
    I(\alpha) = \int_0^1 f_\alpha(x)\dx + \int_1^{+\infty} f_\alpha(x)\dx = I_1(\alpha) + I_2(\alpha)
\]

\paragraph{1. Studio della Singolarità a $x \to 0^+$ ($I_1(\alpha)$):}
Per $x \to 0^+$, si ha $\sqrt{1 + x} \to 1$.
\begin{itemize}
    \item \textbf{Se $\alpha > 0$:} $x^\alpha \to 0$, quindi $\ln(1 + x^\alpha) \sim x^\alpha$.
    Dunque:
    \[
        f_\alpha(x) \sim \frac{x^\alpha}{x^2 \cdot 1} = \frac{1}{x^{2 - \alpha}}
    \]
    Per il criterio del confronto asintotico, $I_1(\alpha)$ converge se e solo se:
    \[
        2 - \alpha < 1 \iff \alpha > 1
    \]
    \item \textbf{Se $\alpha = 0$:} $\ln(1 + x^0) = \ln 2 \implies f_0(x) \sim \frac{\ln 2}{x^2}$, che diverge in $0$ ($2 \ge 1$).
    \item \textbf{Se $\alpha < 0$:} poniamo $\beta = -\alpha > 0$. Per $x \to 0^+$, $x^\alpha = x^{-\beta} \to +\infty$, dunque:
    \[
        \ln(1 + x^{-\beta}) = \ln(x^{-\beta}(1 + x^\beta)) = \beta\ln(1/x) + \bigO(x^\beta) \sim |\alpha|\ln(1/x)
    \]
    L'integrando soddisfa $f_\alpha(x) \sim \frac{|\alpha|\ln(1/x)}{x^2}$. Poiché $\frac{\ln(1/x)}{x^2} \gg \frac{1}{x^{3/2}}$ per $x \to 0^+$, l'integrale diverge.
\end{itemize}
Dunque $I_1(\alpha)$ converge se e solo se $\alpha > 1$.

\paragraph{2. Studio del Comportamento a $x \to +\infty$ ($I_2(\alpha)$):}
Per $x \to +\infty$, $\sqrt{1 + x} \sim x^{1/2}$.
\begin{itemize}
    \item \textbf{Se $\alpha > 0$:} $x^\alpha \to +\infty \implies \ln(1 + x^\alpha) = \ln(x^\alpha(1 + x^{-\alpha})) \sim \alpha\ln x$.
    Dunque:
    \[
        f_\alpha(x) \sim \frac{\alpha\ln x}{x^2 \cdot x^{1/2}} = \frac{\alpha\ln x}{x^{5/2}}
    \]
    Poiché l'esponente $5/2 > 1$, per la scala degli integrali impropri di Bertrand l'integrale $I_2(\alpha)$ converge per ogni $\alpha > 0$.
    \item \textbf{Se $\alpha = 0$:} $f_0(x) \sim \frac{\ln 2}{x^{5/2}}$, convergente a $+\infty$.
    \item \textbf{Se $\alpha < 0$:} $x^\alpha \to 0 \implies \ln(1 + x^\alpha) \sim x^\alpha = \frac{1}{x^{|\alpha|}}$.
    Dunque $f_\alpha(x) \sim \frac{1}{x^{5/2 + |\alpha|}}$, convergente a $+\infty$ poiché $5/2 + |\alpha| > 5/2 > 1$.
\end{itemize}
Pertanto $I_2(\alpha)$ converge per \textbf{ogni} $\alpha \in \R$.

\paragraph{Conclusione Globale:}
L'integrale $I(\alpha)$ converge se e solo se convergono simultaneamente sia $I_1(\alpha)$ che $I_2(\alpha)$:
\[
    \text{Convergenza di } I(\alpha) \iff \alpha > 1
\]
\end{esempio}

\begin{esempio}{Decomposizione in Fratti Semplici e Integrazione Razionale (Appello 2023/24)}{esame_fratti_semplici}
Calcolare il valore esatto dell'integrale definito:
\[
    J := \int_0^1 \frac{x^2 + 2x + 3}{(x+1)(x^2 + 1)} \dx
\]
\tcblower
\textbf{Soluzione Dettagliata:}

\paragraph{1. Decomposizione in Fratti Semplici:}
Il grado del numeratore (2) è strettamente minore del grado del denominatore (3).
Il polinomio al denominatore possiede la radice reale semplice $x = -1$ e il fattore quadratico irriducibile $x^2 + 1$.
Impostiamo la decomposizione:
\[
    \frac{x^2 + 2x + 3}{(x+1)(x^2 + 1)} = \frac{A}{x+1} + \frac{Bx + C}{x^2 + 1}
\]
Moltiplicando per il denominatore comune:
\[
    x^2 + 2x + 3 = A(x^2 + 1) + (Bx + C)(x + 1) = (A + B)x^2 + (B + C)x + (A + C)
\]
Uguagliando i coefficienti dei polinomi (o calcolando nei punti notevoli):
\begin{itemize}
    \item In $x = -1$: $(-1)^2 + 2(-1) + 3 = A(1 + 1) + 0 \implies 2 = 2A \implies A = 1$.
    \item In $x = 0$: $3 = A + C \implies 3 = 1 + C \implies C = 2$.
    \item Coefficiente di $x^2$: $A + B = 1 \implies 1 + B = 1 \implies B = 0$.
\end{itemize}
La decomposizione esatta è quindi:
\[
    \frac{x^2 + 2x + 3}{(x+1)(x^2 + 1)} = \frac{1}{x+1} + \frac{2}{x^2 + 1}
\]

\paragraph{2. Calcolo della Primitiva e Integrazione:}
Integrando termine a termine:
\[
    \int \left( \frac{1}{x+1} + \frac{2}{x^2+1} \right) \dx = \ln|x+1| + 2\arctg(x) + C
\]
Applicando il Teorema Fondamentale del Calcolo Integrale su $[0, 1]$:
\begin{align*}
    J &= \Big[ \ln(x+1) + 2\arctg(x) \Big]_0^1 = \left(\ln(2) + 2\arctg(1)\right) - \left(\ln(1) + 2\arctg(0)\right) \\
    &= \ln 2 + 2\left(\frac{\pi}{4}\right) - 0 = \ln 2 + \frac{\pi}{2}
\end{align*}
\end{esempio}

\section{Equazioni Differenziali Ordinarie (ODE)}
\label{sec:esami_ode}

\begin{esame}[Guida Metodologica: Equazioni Differenziali Ordinarie]
Nelle prove d'esame del Prof. Alberti sulle ODE:
\begin{enumerate}
    \item \textbf{Intervallo Massimale di Esistenza $I_{\max}$ e Fenomeno del Blow-up:}
    Per un problema di Cauchy del I ordine $x'(t) = f(t, x)$ con $f \in C^1$, il Teorema di Cauchy-Lipschitz garantisce esistenza e unicità locale. L'intervallo massimale $I_{\max} = (T_{\min}, T_{\max})$ si interrompe prima di $\pm\infty$ ($T_{\max} < +\infty$) solo se:
    \begin{itemize}
        \item La soluzione diverge in modulo: $\lim_{t \to T_{\max}^-} |x(t)| = +\infty$ (\emph{blow-up} a tempo finito).
        \item La traiettoria $(t, x(t))$ esce dall'aperto di definizione del campo vettoriale (tocca singolarità del tipo $x = 0$ o denominatori nulli).
    \end{itemize}
    Se il campo ha crescita sub-lineare in $x$ ($|f(t, x)| \le A(t)|x| + B(t)$), allora $I_{\max} = \R$ (esistenza globale).
    \item \textbf{Equazioni del II Ordine con Risonanza:}
    Se il termine forzante è $P(t)e^{\alpha t}\cos(\beta t)$ o $P(t)e^{\alpha t}\sen(\beta t)$ e $\lambda = \alpha + i\beta$ è radice del polinomio caratteristico con molteplicità $\mu \ge 1$, la soluzione particolare va cercata con l'ansatz moltiplicato per $t^\mu$.
\end{enumerate}
\end{esame}

\begin{esempio}{Problema di Cauchy a Variabili Separabili ed Esistenza Globale (Appello 2024/25)}{esame_ode_cauchy_blowup}
Risolvere il problema di Cauchy:
\[
    \begin{cases}
        x'(t) = \dfrac{x(t)^2 + 1}{2x(t)} \cdot t \\
        x(0) = \sqrt{3}
    \end{cases}
\]
Determinare la soluzione analitica esplicita $x(t)$, l'intervallo massimale di definizione $I_{\max}$ e il comportamento asintotico per $t \to \pm\infty$.
\tcblower
\textbf{Soluzione Dettagliata:}

\paragraph{1. Separazione delle Variabili e Integrazione:}
L'equazione è a variabili separabili nella forma $x' = g(x)h(t)$ con $g(x) = \frac{x^2+1}{2x}$ e $h(t) = t$.
Il secondo membro è continuo e localmente lipschitziano in $(t, x) \in \R \times (\R \setminus \{0\})$.
Dato che la condizione iniziale è $x(0) = \sqrt{3} > 0$, la soluzione è confinata nel semipiano $x > 0$.
Separando le variabili:
\[
    \frac{2x}{x^2 + 1} \dx = t \dt
\]
Integrando ambo i membri:
\[
    \int \frac{2x}{x^2 + 1} \dx = \int t \dt \implies \ln(x^2 + 1) = \frac{t^2}{2} + C \quad (C \in \R)
\]
Esponenziamo entrambi i membri:
\[
    x^2 + 1 = e^{\frac{t^2}{2} + C} = K e^{t^2/2} \quad (K = e^C > 0)
\]

\paragraph{2. Imposizione del Dato Iniziale:}
Imponendo $x(0) = \sqrt{3}$:
\[
    (\sqrt{3})^2 + 1 = K e^0 \implies 3 + 1 = K \implies K = 4
\]
Dunque:
\[
    x^2(t) + 1 = 4e^{t^2/2} \implies x^2(t) = 4e^{t^2/2} - 1
\]
Poiché $x(0) = \sqrt{3} > 0$, selezioniamo il ramo positivo:
\[
    x(t) = \sqrt{4e^{t^2/2} - 1}
\]

\paragraph{3. Determinazione dell'Intervallo Massimale $I_{\max}$:}
La soluzione è definita e differenziabile fintanto che l'argomento sotto radice è strettamente positivo:
\[
    4e^{t^2/2} - 1 > 0 \iff e^{t^2/2} > \frac{1}{4}
\]
Poiché $e^{t^2/2} \ge 1 > \frac{1}{4}$ per ogni $t \in \R$, la condizione è soddisfatta identicamente per tutti i $t \in \R$.
Inoltre, la soluzione $x(t) \ge \sqrt{4(1) - 1} = \sqrt{3} > 0$ non si avvicina mai alla retta singolare $x = 0$.
Pertanto:
\[
    I_{\max} = (-\infty, +\infty) = \R
\]
La soluzione è globale, è pari ($x(-t) = x(t)$) e possiede il minimo globale in $t = 0$ con $x(0) = \sqrt{3}$.
Per $t \to \pm \infty$:
\[
    \lim_{t \to \pm \infty} x(t) = \lim_{t \to \pm \infty} \sqrt{4e^{t^2/2} - 1} = +\infty
\]
\end{esempio}

\begin{esempio}{Equazione del II Ordine con Risonanza Multipla (Compitino 2024/25)}{esame_ode_2ord_risonanza}
Determinare l'integrale generale dell'equazione differenziale del secondo ordine:
\[
    x''(t) - 4x'(t) + 4x(t) = 6t e^{2t}
\]
e trovare la soluzione che soddisfa le condizioni di Cauchy $x(0) = 1$, $x'(0) = 0$.
\tcblower
\textbf{Soluzione Dettagliata:}

\paragraph{1. Risoluzione dell'Omogenea Associata:}
L'equazione caratteristica associata è:
\[
    P(\lambda) = \lambda^2 - 4\lambda + 4 = (\lambda - 2)^2 = 0
\]
Si ha una radice reale doppia $\lambda_0 = 2$ con molteplicità algebrica $\mu = 2$.
L'integrale generale dell'omogenea è:
\[
    x_{\text{om}}(t) = (c_1 + c_2 t)e^{2t}, \qquad c_1, c_2 \in \R
\]

\paragraph{2. Ricerca della Soluzione Particolare:}
Il termine noto è $f(t) = (6t)e^{2t}$, con esponente $\alpha = 2$ e polinomio $P_1(t) = 6t$ di grado 1.
Essendo $\alpha = 2$ radice caratteristica di molteplicità $\mu = 2$, l'ansatz per la soluzione particolare $\bar{x}(t)$ richiede la moltiplicazione per $t^\mu = t^2$:
\[
    \bar{x}(t) = t^2 (At + B) e^{2t} = (At^3 + Bt^2)e^{2t}
\]
\textit{Metodo rapido dell'operatore shift:}
Notiamo che $\mathcal{L} = (D - 2)^2$. Per la regola di traslazione esponenziale $(D - 2)^2 (u(t)e^{2t}) = u''(t)e^{2t}$.
Posto $u(t) = At^3 + Bt^2$, si ha:
\[
    u'(t) = 3At^2 + 2Bt, \qquad u''(t) = 6At + 2B
\]
Dunque:
\[
    \mathcal{L}(\bar{x}) = u''(t)e^{2t} = (6At + 2B)e^{2t}
\]
Uguagliando al termine forzante $6te^{2t}$:
\[
    (6At + 2B)e^{2t} = 6te^{2t} \iff \begin{cases} 6A = 6 \\ 2B = 0 \end{cases} \iff \begin{cases} A = 1 \\ B = 0 \end{cases}
\]
La soluzione particolare è:
\[
    \bar{x}(t) = t^3 e^{2t}
\]

\paragraph{3. Integrale Generale e Condizioni di Cauchy:}
L'integrale generale è la somma $x(t) = x_{\text{om}}(t) + \bar{x}(t)$:
\[
    x(t) = (c_1 + c_2 t + t^3)e^{2t}
\]
Calcoliamo la derivata prima:
\[
    x'(t) = (c_2 + 3t^2)e^{2t} + 2(c_1 + c_2 t + t^3)e^{2t} = \left[ (2c_1 + c_2) + (2c_2)t + 3t^2 + 2t^3 \right]e^{2t}
\]
Imponendo le condizioni iniziali in $t = 0$:
\[
    \begin{cases}
        x(0) = c_1 = 1 \\
        x'(0) = 2c_1 + c_2 = 0 \implies 2(1) + c_2 = 0 \implies c_2 = -2
    \end{cases}
\]
La soluzione unica del problema di Cauchy è:
\[
    x(t) = (1 - 2t + t^3)e^{2t}
\]
\end{esempio}

\begin{esempio}{ODE Lineare del I Ordine con Parametro e Asintotica (Appello 2023/24)}{esame_ode_lineare_param}
Risolvere il problema di Cauchy:
\[
    \begin{cases}
        y'(t) + \dfrac{2t}{1 + t^2} y(t) = \dfrac{1}{1 + t^2} \\
        y(0) = y_0 \in \R
    \end{cases}
\]
Determinare l'intervallo massimale di esistenza $I_{\max}$ e studiare il limite $\lim_{t \to +\infty} y(t)$ al variare di $y_0$.
\tcblower
\textbf{Soluzione Dettagliata:}

\paragraph{1. Risoluzione con Fattore Integrante:}
L'equazione differenziale è lineare del primo ordine nella forma normale $y'(t) + a(t)y(t) = b(t)$ con $a(t) = \frac{2t}{1+t^2}$ e $b(t) = \frac{1}{1+t^2}$.
Una primitiva del coefficiente $a(t)$ è:
\[
    A(t) = \int_0^t \frac{2s}{1 + s^2}\ds = \ln(1 + t^2)
\]
Il fattore integrante naturale è:
\[
    \mu(t) = e^{A(t)} = e^{\ln(1 + t^2)} = 1 + t^2
\]
Moltiplicando l'equazione per $\mu(t) = 1 + t^2$:
\[
    (1 + t^2)y'(t) + 2t y(t) = 1 \iff \der{}{t}\Big[ (1 + t^2)y(t) \Big] = 1
\]
Integrando tra $0$ e $t$:
\[
    (1 + t^2)y(t) - (1 + 0^2)y(0) = \int_0^t 1\ds = t \implies (1 + t^2)y(t) = y_0 + t
\]
Esplicitando la soluzione:
\[
    y(t) = \frac{y_0 + t}{1 + t^2}
\]

\paragraph{2. Intervallo Massimale e Comportamento Asintotico:}
Poiché il denominatore $1 + t^2 \ge 1 > 0$ non si annulla mai in $\R$, la soluzione è di classe $C^\infty$ su tutto $\R$:
\[
    I_{\max} = (-\infty, +\infty) = \R
\]
Calcoliamo il comportamento asintotico per $t \to +\infty$:
\[
    \lim_{t \to +\infty} y(t) = \lim_{t \to +\infty} \frac{y_0 + t}{1 + t^2} = \lim_{t \to +\infty} \frac{t(1 + y_0/t)}{t^2(1 + 1/t^2)} = \lim_{t \to +\infty} \frac{1}{t} = 0
\]
Il limite vale $0$ per \textbf{qualsiasi} valore della condizione iniziale $y_0 \in \R$.
\end{esempio}

\section{Quesiti Concettuali Rapidi e Test di Parte A}
\label{sec:esami_parteA}

\begin{esame}[Strategie per il Test Preliminare di Parte A]
Il test di Parte A (quesiti a risposta sintetica o multipla) richiede rapidità e precisione geometrico-analitica:
\begin{enumerate}
    \item \textbf{Dominio e Immagine:} Per calcolare $\Imm(f)$, determinare prima gli intervalli di monotonia della funzione nei vari rami del dominio e calcolare i limiti agli estremi per ciascun ramo continuo.
    \item \textbf{Coordinate Polari:} Ricordare la determinazione dell'angolo polare $\theta \in [0, 2\pi)$:
    \[
        \theta = \begin{cases}
            \arctg(y/x) & \text{nel I quadrante } (x > 0, y \ge 0) \\
            \arctg(y/x) + \pi & \text{nel II e III quadrante } (x < 0) \\
            \arctg(y/x) + 2\pi & \text{nel IV quadrante } (x > 0, y < 0)
        \end{cases}
    \]
    \item \textbf{Estremo Superiore/Inferiore di Insiemi Discreti:} Separare sempre le sottosuccessioni a termini di indice pari e dispari per identificare gli andamenti monotoni verso gli estremi.
\end{enumerate}
\end{esame}

\begin{esempio}{Insieme di Definizione, Immagine e Inversa (Compitino 2024/25)}{esame_parteA_inversa}
Determinare l'insieme di definizione, l'immagine e l'inversa della funzione reale:
\[
    f(x) := 2 - \exp\left(\frac{1}{x+1}\right)
\]
\tcblower
\textbf{Soluzione Dettagliata:}
\begin{enumerate}
    \item \textbf{Dominio naturale}: la funzione esponenziale è definita ovunque, mentre l'esponente richiede la non nullità del denominatore $x+1 \neq 0$. Dunque:
    \[
        \dom(f) = \R \setminus \{-1\} = (-\infty, -1) \cup (-1, +\infty)
    \]
    \item \textbf{Immagine}: studiamo il comportamento sui due intervalli connessi:
    \begin{itemize}
        \item Per $x \in (-1, +\infty)$: la funzione $u(x) = \frac{1}{x+1}$ è strettamente decrescente da $+\infty$ a $0^+$.
        Di conseguenza $e^{u(x)} \in (1, +\infty)$, da cui:
        \[
            y = 2 - e^{u(x)} \in (-\infty, 1)
        \]
        \item Per $x \in (-\infty, -1)$: la funzione $u(x) = \frac{1}{x+1}$ è strettamente decrescente da $0^-$ a $-\infty$.
        Di conseguenza $e^{u(x)} \in (0, 1)$, da cui:
        \[
            y = 2 - e^{u(x)} \in (1, 2)
        \]
    \end{itemize}
    L'immagine complessiva è l'unione dei due intervalli:
    \[
        \Imm(f) = (-\infty, 1) \cup (1, 2) = \{y \in \R \mid y < 2 \text{ e } y \neq 1\}
    \]
    \item \textbf{Funzione Inversa}: poiché $f$ è strettamente monotona su ciascun intervallo e le due immagini sono disgiunte, $f$ è globale iniettiva su $\dom(f)$. Esplicitiamo $x$ in funzione di $y \in \Imm(f)$:
    \[
        y = 2 - e^{\frac{1}{x+1}} \iff e^{\frac{1}{x+1}} = 2 - y \iff \frac{1}{x+1} = \ln(2 - y) \iff x + 1 = \frac{1}{\ln(2 - y)}
    \]
    Dunque la funzione inversa $f^{-1} \colon (-\infty, 1) \cup (1, 2) \to \R \setminus \{-1\}$ è:
    \[
        f^{-1}(y) = \frac{1}{\ln(2 - y)} - 1
    \]
\end{enumerate}
\end{esempio}

\begin{esempio}{Coordinate Polari nel Piano Reale (Compitino 2024/25)}{esame_parteA_polari}
Determinare le coordinate polari $(\rho, \theta)$ con raggio $\rho \ge 0$ e anomalia $\theta \in [0, 2\pi)$ dei seguenti punti del piano cartesiano:
\[
    P_1(-1, \sqrt{3}), \qquad P_2(0, -2), \qquad P_3(3, -3)
\]
\tcblower
\textbf{Soluzione Dettagliata:}
Ricordando che $\rho = \sqrt{x^2 + y^2}$ e $\theta$ tiene conto dell'angolo di rotazione antioraria a partire dal semiasse positivo delle ascisse:
\begin{enumerate}
    \item \textbf{Punto $P_1(-1, \sqrt{3})$:} appartiene al \textbf{II quadrante} ($x < 0, y > 0$).
    \[
        \rho_1 = \sqrt{(-1)^2 + (\sqrt{3})^2} = \sqrt{1 + 3} = 2
    \]
    \[
        \theta_1 = \arctg\left(\frac{\sqrt{3}}{-1}\right) + \pi = -\frac{\pi}{3} + \pi = \frac{2\pi}{3}
    \]
    \item \textbf{Punto $P_2(0, -2)$:} giace sul semiasse negativo delle ordinate ($x = 0, y < 0$).
    \[
        \rho_2 = \sqrt{0^2 + (-2)^2} = 2, \qquad \theta_2 = \frac{3\pi}{2}
    \]
    \item \textbf{Punto $P_3(3, -3)$:} appartiene al \textbf{IV quadrante} ($x > 0, y < 0$).
    \[
        \rho_3 = \sqrt{3^2 + (-3)^2} = \sqrt{9 + 9} = 3\sqrt{2}
    \]
    \[
        \theta_3 = \arctg\left(\frac{-3}{3}\right) + 2\pi = \arctg(-1) + 2\pi = -\frac{\pi}{4} + 2\pi = \frac{7\pi}{4}
    \]
\end{enumerate}
\end{esempio}

\begin{esempio}{Cinematica, Velocità Istantanea e Lunghezza di Curva (Compitino 2024/25)}{esame_parteA_curva}
Sia dato un punto materiale $P$ che si muove nello spazio tridimensionale con traiettoria parametrizzata dalla legge oraria:
\[
    P(t) := \left( e^t + e^{-t}, \, \cos(2t), \, \sen(2t) \right) \in \R^3, \quad t \in \R
\]
Calcolare il modulo del vettore velocità istantanea $|\vec{v}(t)|$ e la lunghezza complessiva dell'arco di curva descritto nell'intervallo $t \in [-2, 2]$.
\tcblower
\textbf{Soluzione Dettagliata:}
Il vettore velocità istantanea $\vec{v}(t) = P'(t)$ è ottenuto derivando componente per componente:
\[
    \vec{v}(t) = \left( e^t - e^{-t}, \, -2\sen(2t), \, 2\cos(2t) \right)
\]
Calcoliamo la velocità scalare istantanea (modulo del vettore velocità):
\begin{align*}
    |\vec{v}(t)| &= \sqrt{(e^t - e^{-t})^2 + (-2\sen(2t))^2 + (2\cos(2t))^2} \\
    &= \sqrt{e^{2t} - 2 + e^{-2t} + 4\left(\sen^2(2t) + \cos^2(2t)\right)} \\
    &= \sqrt{e^{2t} - 2 + e^{-2t} + 4} = \sqrt{e^{2t} + 2 + e^{-2t}} = \sqrt{(e^t + e^{-t})^2} = e^t + e^{-t}
\end{align*}
La lunghezza $L$ dell'arco di curva nell'intervallo $[-2, 2]$ è l'integrale del modulo della velocità:
\[
    L = \int_{-2}^2 |\vec{v}(t)|\dt = \int_{-2}^2 (e^t + e^{-t})\dt = \Big[ e^t - e^{-t} \Big]_{-2}^2 = (e^2 - e^{-2}) - (e^{-2} - e^2) = 2(e^2 - e^{-2})
\]
\end{esempio}

\begin{esempio}{Estremo Superiore, Inferiore, Massimo e Minimo di Insiemi Numerici (Compitino 2023/24)}{esame_parteA_sup_inf}
Determinare l'estremo superiore, l'estremo inferiore e l'eventuale esistenza di massimo e minimo per l'insieme numerico:
\[
    E := \left\{ \frac{2n (-1)^n}{n + 1} \in \R \;\middle|\; n \in \N, \, n \ge 1 \right\}
\]
\tcblower
\textbf{Soluzione Dettagliata:}
Analizziamo gli elementi $a_n = \frac{2n(-1)^n}{n+1}$ distinguendo la parità dell'indice $n$:
\begin{enumerate}
    \item \textbf{Indici pari ($n = 2k$ con $k \ge 1$):}
    \[
        a_{2k} = \frac{2(2k)(-1)^{2k}}{2k + 1} = \frac{4k}{2k + 1} = 2 - \frac{2}{2k + 1} > 0
    \]
    La successione dei termini pari $a_{2k}$ è strettamente crescente rispetto a $k$:
    \begin{itemize}
        \item Per $k = 1$ ($n = 2$): $a_2 = \frac{4}{3}$.
        \item Per $k \to +\infty$: $\lim_{k \to \infty} a_{2k} = 2$.
    \end{itemize}
    Poiché $a_{2k} < 2$ per ogni $k \ge 1$, si ha:
    \[
        \sup E = 2, \qquad \max E \text{ non esiste (poiché } 2 \notin E\text{)}
    \]
    \item \textbf{Indici dispari ($n = 2k - 1$ con $k \ge 1$):}
    \[
        a_{2k-1} = \frac{2(2k-1)(-1)^{2k-1}}{2k} = -\frac{2k - 1}{k} = -2 + \frac{1}{k} < 0
    \]
    La successione dei termini dispari $a_{2k-1}$ è strettamente decrescente rispetto a $k$:
    \begin{itemize}
        \item Per $k = 1$ ($n = 1$): $a_1 = -2 + 1 = -1$.
        \item Per $k \to +\infty$: $\lim_{k \to \infty} a_{2k-1} = -2$.
    \end{itemize}
    Poiché $a_{2k-1} > -2$ per ogni $k \ge 1$, si ha:
    \[
        \inf E = -2, \qquad \min E \text{ non esiste (poiché } -2 \notin E\text{)}
    \]
\end{enumerate}
\paragraph{Riepilogo:}
\[
    \sup E = 2 \quad (\nexists \max E), \qquad \inf E = -2 \quad (\nexists \min E)
\]
\end{esempio}

